\begin{afterword}
\summaryhead

The post argues that people expect others to understand their words as they intend them, because they themselves know what they mean. It calls this the illusion of transparency and relates it to hindsight bias.

It reports several experiments by Boaz Keysar and colleagues. Subjects who knew that a restaurant was bad tended to predict that a listener would hear sarcasm in a message praising it. Subjects taught one meaning of an obscure idiom expected uninformed listeners to arrive at the same meaning. Speakers who said ambiguous sentences estimated that they were understood 72\% of the time, when they were understood 61\% of the time; listeners who overheard, knowing the intended meaning, did not show this overestimate.

The post adds a historical case, a letter from Chamberlain that it says Hitler read as conciliatory before invading Poland. It concludes that you should not be quick to blame people who misunderstand you, because your words are probably more ambiguous than you think.

\responsehead

The post has one real finding to report: speakers overestimate how often they are understood. It reports that finding in a way that commits the error the finding describes.

Start with the title. By the post's own figures, listeners understood the speakers 61\% of the time, and the speakers expected 72\%. That is an overestimate of eleven points. The title turns it into ``No One Understands You.'' If the author expected readers to take ``no one'' as rhetoric and not as a claim, he was doing what the speakers in the study did: trusting his intention to come through the words. The title also does something else. It speaks to a reader who feels misunderstood and tells them why. The body then asks that reader to be more careful with their own words. The title is what most people remember.

Next, the mechanism. The post says we know what we mean, ``and so we expect others to know it too.'' The Keysar and Henly study included overhearers who were told the intended meaning. They knew what the speaker meant, and they did not systematically overestimate. So knowing the meaning is not enough to produce the bias. The study's authors locate it in speakers monitoring ``their own utterances.'' The post reports the overhearer result in one sentence, as a separate paragraph, and does not see that it undercuts the paragraph that defined the illusion.

Then the sourcing. Two of the three findings in the sarcasm paragraph come through a secondary chapter. The post guesses whether a voice recording was used in a paper whose abstract says subjects ``read scenarios.'' It calls June ``Jane'' twice. It gives the editors of its main secondary source as ``Griffin Gilovich and Daniel Kahneman''; the editors are Thomas Gilovich, Dale Griffin and Daniel Kahneman. It uses the name of a 1998 paper by Gilovich and colleagues without crediting it.

The Chamberlain story, relayed at second hand, is wrong on the date, the wording and the reading. The letter was sent ten days before the invasion, not two. It is not phrased in vague diplomatese; it says Britain is ``resolved, and prepared, to employ without delay all the forces at their command.'' Hitler's reply shows that he understood: ``I take note of this statement.'' A post about the gap between what a writer means and what a reader takes away closes its evidence with a letter it did not check.

In short: a post on overconfidence in one's own words, with a title its own data refute, a mechanism its own control group undercuts, and a closing example it never checked.
\end{afterword}
